%%% SECTION 7: PHYSICAL AI STANDARD LEVEL AND UNIFICATION STANDARD %%%

The Physical AI Standard Level (PSL) and Unification Standard Level (USL)
frameworks complement each other to create a comprehensive, quantitative
standards system for responsible robot usage in oncology clinical trials. PSL
provides the regulatory compliance dimensions for trial site operations, while
USL evaluates individual robot readiness. Together, they ensure that both the
site infrastructure and the robots operating within it meet the standards
required for safe, effective Physical AI oncology trials. This dual-layer
quantitative assessment replaces subjective evaluation with reproducible,
objective scoring that enables standardized nationwide deployment.

\subsection{Overview: Complementary Standards}
\label{subsec:standards-overview}

PSL and USL serve different but complementary purposes. PSL evaluates
site-level regulatory compliance across three dimensions, determining whether a
trial site has met the legislative, regulatory, and operational requirements
for Physical AI trial operations. USL evaluates individual robot readiness
across four dimensions, determining whether a specific robot platform is
sufficiently mature for clinical trial deployment. A trial site can only
operate when both conditions are satisfied: the site must achieve passing PSL
scores (site is compliant) and the robots deployed at that site must achieve
sufficient USL scores (robots are ready).

This dual-layer approach is unique to this National Platform and more thorough
than any existing FDA guidance on robot or AI evaluation. Existing readiness
frameworks such as NASA Technology Readiness Levels (TRL 1-9) evaluate hardware
maturity but do not address clinical regulatory compliance, and existing AI
evaluation frameworks do not capture the physical embodiment requirements of
robotic systems in clinical settings.

\subsection{Physical AI Standard Level (PSL) Framework}
\label{subsec:psl-framework}

The PSL framework was developed in conjunction with the Clinical Trial Site
Complete Documentation \cite{site-docs} and validated through the 24-hour
clinical trial simulation that treated 168 patients with 29 robots.

\subsubsection{The Three PSL Dimensions}

PSL evaluates site-level compliance across three regulatory dimensions that
correspond to the three layers of authorization required for Physical AI trial
site operation.

\begin{table}[H]
\centering
\caption{Physical AI Standard Level (PSL) Three Dimensions}
\label{tab:psl-dimensions}
\small
\begin{tabularx}{\textwidth}{c l X}
\toprule
\textbf{Dim} & \textbf{Name} & \textbf{Scope} \\
\midrule
1 & Legislative Authorization & State legislation enabling Physical AI trials (SB 1042, AB 2847, SB 892), federal statutory authority, municipal code compliance \\
2 & Regulatory Compliance & Title 22 state regulations, FDA compliance guide, building and premises codes, 21 CFR Part compliance \\
3 & Operational Readiness & Activation SOPs, emergency preparedness, staff training, equipment verification, PSL score verification \\
\bottomrule
\end{tabularx}
\end{table}

Dimension 1 (Legislative Authorization) measures whether the required
legislative framework exists and applies to the trial site. For California,
this includes SB 1042 (Physical AI Oncology Clinical Trial Authorization Act),
AB 2847 (Patient Rights and Robotic Safety Act), and SB 892 (Data Protection
and Transparency Act). Each legislative instrument is scored on completeness,
applicability, and enforcement readiness.

Dimension 2 (Regulatory Compliance) measures whether the site meets all
regulatory requirements across state and federal jurisdictions. This includes
Title 22 regulations for healthcare facility operations, FDA compliance with
the three adapted standards \cite{ich-adapt, cfr50-adapt, cfr312-adapt},
building code standards for robot operations, and premises code requirements
for robot storage, charging, and maintenance areas.

Dimension 3 (Operational Readiness) measures whether the site has completed all
preparatory activities required for safe Physical AI operations. This includes
activation standard operating procedures, emergency preparedness plans, staff
training and certification, equipment verification and calibration, and formal
PSL score verification by an independent assessor.

\subsubsection{PSL Scoring Methodology}

PSL scores are computed on a pass-fail basis for each dimension, with granular
sub-scores for individual requirements within each dimension. A site must
achieve passing status on all three dimensions before activation. Within each
dimension, individual requirements are scored as: Compliant (requirement fully
met with documentation), Partially Compliant (requirement substantially met
with documented plan for completion), or Non-Compliant (requirement not met).
A dimension passes when all critical requirements are Compliant and all
non-critical requirements are at least Partially Compliant.

\subsubsection{PSL and Site Activation}

PSL scores determine whether a trial site can be activated for Physical AI
operations. The activation process follows the procedures detailed in the site
establishment documentation (Section~\ref{sec:site-establishment}), with PSL
score verification serving as a gate before operational launch. The 24-hour
simulation \cite{site-docs} validated this activation process by demonstrating
that a site meeting PSL requirements could safely operate 29 robots serving 168
patients continuously.

\subsection{Unification Standard Level (USL) Framework}
\label{subsec:usl-framework}

The USL framework \cite{usl-standard} provides a 1.0 to 10.0 scoring system
for evaluating robot readiness across four equally weighted dimensions. USL was
developed collaboratively using Claude Code Opus 4.6 \cite{claude-opus} with
scoring engines implemented in Python across three category-specific modules
(cobots, surgical robots, humanoid robots) totaling over 2,600 lines of code.

\subsubsection{The Four USL Dimensions}

\begin{table}[H]
\centering
\caption{Unification Standard Level (USL) Four Dimensions}
\label{tab:usl-dimensions}
\small
\begin{tabularx}{\textwidth}{c l c X}
\toprule
\textbf{Dim} & \textbf{Name} & \textbf{Wt} & \textbf{What It Measures} \\
\midrule
A & Simulation Switching & 25\% & Transfer of policies between simulation engines (Isaac Lab, MuJoCo, Gazebo, PyBullet, Drake) \\
B & AI Integration & 25\% & Integration with LLMs, VLAs, diffusion policies, MCP, agentic AI frameworks \\
C & Cross-Robot Sharing & 25\% & Intra- and inter-organization progress sharing via ONNX, ROS 2, federated learning \\
D & Clinical Trial Collab & 25\% & Regulatory compliance, federated deployment, audit trails across clinical trial sites \\
\bottomrule
\end{tabularx}
\end{table}

Dimension A (Simulation Framework Switching) evaluates a robot's ability to
transfer trained policies between simulation engines. Robots that operate in
only one simulation environment score lower than those with validated
multi-engine support. This dimension captures the practical reality that
different trial sites may use different simulation platforms, and robots must
function reliably across these environments.

Dimension B (AI Integration) evaluates the depth and breadth of AI system
integration. This encompasses large language model planning capabilities,
vision-language-action (VLA) model support, diffusion policy compatibility,
Model Context Protocol \cite{mcp-protocol} support, and agentic AI framework
integration. Higher scores indicate robots that can leverage the full spectrum
of modern AI capabilities.

Dimension C (Cross-Robot Sharing) evaluates the ability to share learned
behaviors, trained models, and operational data across robots within and
between organizations. ONNX model export, ROS 2 action interfaces, checkpoint
sharing, and federated learning compatibility are key scoring factors. This
dimension is essential for scaling Physical AI across multiple trial sites.

Dimension D (Clinical Trial Collaboration) evaluates regulatory compliance
status, clinical deployment history, federated learning readiness, audit trail
capabilities, and multi-site deployment infrastructure. This is the most
directly relevant dimension for Physical AI oncology trials and was found to
be the weakest dimension for seven of nine robots evaluated.

\subsubsection{USL Score Computation}

The final USL score is computed as the weighted average of all four dimension
scores, clamped to the range 1.0 to 10.0 and rounded to the nearest 0.1:

$$\text{USL}_{\text{final}} = \text{clamp}\left(\frac{0.25 \cdot S_A + 0.25 \cdot S_B + 0.25 \cdot S_C + 0.25 \cdot S_D}{1.0},\; 1.0,\; 10.0\right)$$

Scores map to five bands: Exemplary (9.0-10.0), Advanced (7.0-8.9),
Intermediate (5.0-6.9), Foundational (3.0-4.9), and Initial (1.0-2.9). Each
band corresponds to a set of qualitative descriptors for clinical trial
readiness, from fully unified multi-site deployment capability (Exemplary) to
minimal unification capability (Initial).

\subsection{USL Results Across Robot Categories}
\label{subsec:usl-results}

Table~\ref{tab:usl-all-scores} presents the complete USL evaluation results for
all nine robots across three categories, providing the quantitative foundation
for robot selection in Physical AI oncology trials.

\begin{table}[H]
\centering
\caption{Complete USL Scores for All Nine Evaluated Robots}
\label{tab:usl-all-scores}
\small
\begin{tabular}{l l r r r r r l}
\toprule
\textbf{Robot} & \textbf{Manufacturer} & \textbf{A} & \textbf{B} & \textbf{C} & \textbf{D} & \textbf{USL} & \textbf{Band} \\
\midrule
\multicolumn{8}{l}{\textit{Collaborative Robots (Cobots)}} \\
Franka Panda & Franka Robotics & 8.0 & 7.7 & 8.5 & 5.5 & \textbf{7.4} & Advanced \\
Kinova Gen3 & Kinova Robotics & 6.3 & 5.2 & 5.8 & 5.5 & \textbf{5.7} & Intermediate \\
xArm 7 & UFACTORY & 4.6 & 3.2 & 3.8 & 2.1 & \textbf{3.4} & Foundational \\
\midrule
\multicolumn{8}{l}{\textit{Surgical Robots}} \\
da Vinci dVRK & Intuitive Surgical & 7.5 & 7.2 & 6.8 & 7.0 & \textbf{7.1} & Advanced \\
Hugo RAS & Medtronic & 4.5 & 4.0 & 3.5 & 5.8 & \textbf{4.5} & Foundational \\
Versius & CMR Surgical & 3.2 & 2.8 & 2.5 & 5.2 & \textbf{3.4} & Foundational \\
\midrule
\multicolumn{8}{l}{\textit{Humanoid Robots}} \\
Atlas Electric & Boston Dynamics & 7.0 & 7.4 & 4.5 & 4.2 & \textbf{5.8} & Intermediate \\
Digit & Agility Robotics & 5.8 & 5.4 & 3.0 & 2.7 & \textbf{4.2} & Foundational \\
Optimus Gen 2 & Tesla & 3.4 & 5.0 & 1.5 & 4.4 & \textbf{3.6} & Foundational \\
\bottomrule
\end{tabular}
\end{table}

\subsubsection{Surgical Robot Results}

The da Vinci dVRK leads surgical robots with USL 7.1 (Advanced band) and
achieves the highest Dimension D score (7.0) of any robot evaluated. This
reflects FDA clearance, deployment at over 9,000 systems worldwide, 14 million
procedures across 69 countries, and an open-source research platform used by
45 institutions. Franka Panda (USL 7.4) outperforms da Vinci overall due to
superior simulation switching (8.0 vs 7.5) and cross-robot sharing (8.5 vs
6.8), reflecting its broader open-source ecosystem.

Hugo RAS (USL 4.5, Foundational) and Versius (USL 3.4, Foundational) score
lower primarily due to proprietary ecosystems that limit simulation switching,
AI integration, and cross-robot sharing. Both achieve relatively higher
Dimension D scores (5.8 and 5.2 respectively) reflecting CE marking and active
clinical deployment in European markets.

\subsubsection{Cobot and Humanoid Results}

Among cobots, the Franka Panda's dominance (USL 7.4) is driven by support
across five simulation frameworks, five AI capabilities, and four sharing
methods including inter-organizational ONNX transfer. Kinova Gen3 (USL 5.7,
Intermediate) offers strong clinical workflow compatibility but lacks
VLA/diffusion policy support. The xArm 7 (USL 3.4, Foundational) demonstrates
that hardware excellence (IP51 protection, built-in collision detection) does
not compensate for limited software ecosystem maturity.

Among humanoids, Atlas Electric (USL 5.8, Intermediate) achieves the highest AI
integration score (7.4) of any humanoid, driven by foundation model potential
and whole-body control capabilities. Digit (USL 4.2, Foundational) benefits
from the NVIDIA GR00T partnership but lacks healthcare deployment. Optimus
Gen 2 (USL 3.6, Foundational) has the most capable hands (11 DOF, 22 joints)
and mass production potential but scores the lowest cross-robot sharing score
(1.5) of all nine robots due to its fully proprietary platform.

\subsubsection{Cross-Category Comparison}

The cross-category comparison reveals two key findings. First, open-source
ecosystem size is the strongest predictor of USL score. The two highest-scoring
robots (Franka 7.4, da Vinci 7.1) have the largest open-source communities.
Second, Dimension D (Clinical Trial Collaboration) is the weakest dimension for
seven of nine robots, revealing a field-wide gap between research maturity and
clinical deployment infrastructure. Even the highest-scoring robot (Franka,
7.4 overall) scores only 5.5 on Dimension D.

\subsection{How PSL and USL Complement Each Other}
\label{subsec:psl-usl-complement}

The three PSL dimensions and four USL dimensions work together as a dual-layer
assessment. PSL ensures the environment is ready (site compliance); USL ensures
the robots are ready (platform readiness). A Physical AI trial site requires
both: a compliant, well-prepared facility (passing PSL) populated with properly
qualified robots (sufficient USL scores).

A site might pass PSL on all three dimensions (legislation in place, regulatory
compliance achieved, operations prepared) but be unable to deploy certain robot
types because their USL scores are insufficient. For example, a California
trial site meeting all PSL requirements could deploy da Vinci dVRK (USL 7.1,
Advanced) for surgical procedures but might need to defer deployment of xArm 7
(USL 3.4, Foundational) until its software ecosystem matures. Conversely, a
robot with excellent USL scores cannot operate at a site that fails PSL
requirements, regardless of the robot's individual capabilities.

This complementary relationship is connected to the MCP conformance levels from
Physical AI National MCP Servers \cite{national-mcp-paper}, which provide an
additional interoperability standard. MCP conformance levels interface with
both PSL (ensuring site infrastructure supports MCP communication) and USL
(evaluating whether robots support standardized MCP protocols). The Clinical
Trial Site Complete Documentation \cite{site-docs} and Physical AI Unification
Standard Level \cite{usl-standard} serve as the primary reference
implementations for PSL and USL respectively.

\subsection{Impact on Nationwide Deployment}
\label{subsec:standards-impact}

Standardized PSL and USL scores enable efficient nationwide scaling of Physical
AI oncology trials. When a new site seeks to establish Physical AI trial
operations, the PSL framework provides a clear checklist of legislative,
regulatory, and operational requirements that must be met. When a sponsor
evaluates which robots to deploy at a new site, USL scores provide objective
data on each robot's readiness.

The open-source nature of both frameworks, available through the three code
repositories \cite{main-repo, mcp-repo, fl-repo}, ensures that any institution
can access, evaluate, and apply the scoring methodology. This transparency
supports the broader goal of democratizing Physical AI oncology trial access
rather than limiting it to institutions with proprietary evaluation processes.

The field-wide weakness in Dimension D (Clinical Trial Collaboration) identified
by USL scoring highlights a critical area for targeted improvement. As the
Physical AI ecosystem matures, Dimension D scores should improve as robots gain
regulatory clearance, federated learning compatibility, audit trail
capabilities, and multi-site deployment experience. The PSL and USL frameworks
provide the measurement infrastructure to track this progress over time.
